\authoraction{Complete the peripersonal-space results, verify the reporting period and ethics reference, and confirm publication status before submission. This draft is not a declaration that the project has been successfully concluded.}

\section{Aims}

The project investigates whether coupling respiration to changes in an immersive virtual environment can alter conscious experience and the perceived relationship between the body and surrounding space. The approved programme comprised Study I, development of a bioresponsive virtual-reality system, and Study II, evaluation of its effects using subjective measures of altered states and objective behavioural measures of peripersonal space. [1]

We developed Viscereality as a breathing-responsive particle environment and subsequently evaluated a controlled box-breathing implementation. The completed subjective study supports differences between immersive VR and a minimal-feedback control in breathing adherence, spatial self-reference and several altered-state dimensions. It did not show reliable differences between the two tested visual symmetry conditions. The behavioural peripersonal-space component addresses a separate part of the approved aims and must be reported separately. [2--5]

\subsection{Project implementation and approved adjustments}

The original Study II proposed a 2 \(\times\) 2 mixed design combining breathwork versus ordinary breathing with congruent versus incongruent spatial feedback, a target of 30 participants, and both subjective and peripersonal-space outcomes. The implemented programme separated subjective and behavioural investigations to reduce session length and participant burden and to accommodate technical constraints. In the explanation prepared in April 2026, we described the difficulty of combining the two objectives within a single sustained VR session. [1,2]

The implemented programme incorporates methodological adjustments accepted by the Foundation. This draft distinguishes the approved aims, the experiments actually performed and the evidence still required to establish completion. Approval of adjustments does not establish acceptance of the final report or fulfilment of the journal-publication requirement. [2]

{
\small\setlength{\tabcolsep}{4pt}\renewcommand{\arraystretch}{1.16}
\begin{longtable}{>{\raggedright\arraybackslash}p{\dimexpr 0.20585\linewidth-2\tabcolsep\relax}>{\raggedright\arraybackslash}p{\dimexpr 0.36470\linewidth-2\tabcolsep\relax}>{\raggedright\arraybackslash}p{\dimexpr 0.42945\linewidth-2\tabcolsep\relax}}
\toprule
\textbf{Approved element} & \textbf{Implemented work} & \textbf{Interpretive boundary} \\
\midrule\endfirsthead
\toprule
\textbf{Approved element} & \textbf{Implemented work} & \textbf{Interpretive boundary} \\
\midrule\endhead
\bottomrule\endfoot
Study I VR development & Bioresponsive particle system; respiratory and cardiac design outputs & System development is documented by two conference papers. \\
Study II subjective outcomes & Counterbalanced three-condition box-breathing study & This is not the original breathwork \(\times\) congruency factorial experiment. \\
Study II objective spatial outcomes & Separate preregistered peripersonal-space study & Final sample, exclusions, analyses and results must be inserted. \\
Expected publications & Two conference papers and a journal manuscript & Conference papers do not satisfy the required journal article. \\
\end{longtable}
}
